\newpage
\section{Formalization}
% Id = {id1, id2, ...}
% Cst = some constant domain

% Type = 'Id | 'Cst
% Schema = [Type] 

% DB = (ConstFun, IdFun, I, =_Id, R) where
%     ConstFun: Symbol -> (Schema, Cst x Cst -> Cst)
%     IdFun: Symbol -> Schema
%     // Fun: SYmbol -> (Schema, CstxCst->Cst | 'union)
%     I: (s:Symbol) -> Set<D^(arity(s)+1)>
%     R: [Rule]

% Rule ::= [P] => [A]
% P ::= (= Var Pat)
% A ::= (set (func Pat ...) Pat)
% Pat ::= (func Pat ...) | Cst | Var

% (schema, merge) = ConstFun(f)
% DB = (ConstFun, IdFun, I, =_Id, R)
% DB' = (ConstFun, IdFun, I', =_Id, R)
% c_i = eval(Gamma, p_i) for i = 1 ... k
% c =  eval(Gamma, p)
% // TODO: c' may not exist
% c' = { c' : (c1, ..., ck, c') in I(f) }
% I' = I[f -> I(f) - {(c1, ..., ck, c')} + {(c1, ..., ck, merge(c, c')) ]
% ------------------------------------
% Gamma |- DB -[ (set (f p1 ... pk) p) ]-> DB'


% DB' = ICO(DB)

% TODO: if we define the lattice as constants, 
% then you can do "plain" matching over the lattices,
% which is not semantically sound. 
% The bigger question is, should we disallow non-monotonic operations
% Well, the lambda example has it.
% you can syntactically restrict to something like (= c (max (lo ...) (lo ...))), so that it's always sound.
% this is not Zach's "VM view" though...

\def\DB{\mathit{DB}}
\def\bmt{\textbf{t}}
\def\mergef{\mathit{merge}}

Given an infinite set of uninterpreted constants $N = \{n_1,n_2, \ldots\}$ 
 and a bounded semi-lattice of (interpreted) constants $L = (C, \bot, \sqcup)$,
 a schema $S$ is a collection of function symbols and their function signatures, 
 where the types range over $\{N, C\}$.
Given a schema $S$, an instance of $S$ is defined $I=(\DB, \equiv)$,
 where $\DB$ is as a mapping from a function symbol $f$ in $S$ 
 to a partial function implementing the associated signature, 
 and $\equiv$ is an equivalence relation over $N\cup C$ 
 satisfying $\forall c_1, c_2\in C. c_1\equiv c_2\rightarrow c_1=c_2$ (i.e., constants are only equivalent to themselves).

Given an arbitrary total order $<$ over $N\cup C$,
we define the canonicalization function $\lambda_{\equiv}(t)=\min t':t'\equiv t$.
We say $f(\bmt{})\mapsto t\in I$ if $\DB{}[f](\lambda_{\equiv}(\bmt{}){})\equiv t$.
For convenience, we lift $\lambda_{\equiv}$ to 
 also work on sets 
 and the whole database instances by pointwise canonicalization.

Given an \egglog program $P$ 
 consisting of a set of rules and $I=(\DB{}, \equiv)$, 
 we define the immediate consequence operator $T_P(I)=(\DB{}', \equiv)$, where
\begin{align*}
    \DB{}'[f]=&\{ \bmt{}\mapsto t\;|\;f(\bmt{})\mapsto t\;\textit{:-}\;F_1,\ldots,F_n \text{ is an instantiation of a rule in $P$ where $F_i\in I$ } \}\\
    \cup& \{ \bmt{}\mapsto \textit{default}_f\;|\;\exists t.f(\bmt{})\mapsto t\;\textit{:-}\;F_1,\ldots,F_n \text{ is an instantiation of a rule in $P$ where $F_i\in I$}\\
    &\text{\quad\quad\quad\quad\quad\quad\  and $\not \exists t. f(\bmt{})\mapsto t\in I$\} } 
\end{align*}
where $\textit{default}_f$ is a fresh constant from $N$ not occurring in the database if the output type of $f$ is $N$ and is $\bot$ otherwise.

Note however $\DB{}'[f]$ may no longer be functions, 
 as it is possible that the same key $\bmt{}$ can now be mapped to more than one $t$.
We call $T_P(I)$ a ``pre-instance'', since they are not valid instances yet.
To transform a pre-instance into a valid instance, 
 we further define the rebuilding operator \(
    R(I) = (\DB{}_R, \equiv_R)
\), 
where $\equiv_R$ is the equivalence closure of 
\[
    \equiv\cup 
    \{(t_1,t_2)\;|f(\bmt{})\mapsto t_1\in I, 
    f(\bmt{})\mapsto t_2\in I\}
\]
and
\[
    % \DB{}_R[f]=\left\{
    %     \bmt\mapsto 
    %     \mergef_{f,\equiv}\left(
    %         \left\{t: f\left(\bmt\right)\mapsto t\in I\right\}
    %     \right)\;|\; 
    %     \text{for all $\bmt$ where $\left\{t: f\left(\bmt\right)\mapsto t\in I\right\}$ is not empty}
    % \right\}
    \DB{}_R[f]=\lambda_{\equiv_R}\left(\left\{
    \bmt\mapsto 
    \mergef_{f,\equiv}\left(
        K
    \right)\;|\; 
    \text{ $K=\left\{t: f(\bmt)\mapsto t\in I\right\}$ and $K$ is not empty}
    \right\}\right)
\]
% where \[
%     \mergef_{f,\equiv}(K)=
%     \min\left(\lambda_{\equiv}(K)\right)
% \]
% if the output type of $f$ is $N$ and \[
%     \mergef_{f,\equiv}(K)=\bigsqcup K
% \] if the output type of $f$ is $C$.
where \[
    \mergef_{f,\equiv}(K)=
    \begin{cases}
        \min\left(\lambda_{\equiv}(K)\right) &\text{if the output type of $f$ is $N$;}\\
        \bigsqcup K &\text{if the output type of $f$ is $C$.}
    \end{cases}
\]

\def\lfp{\mathit{lfp}}

Note that the canonicalization in computing $\DB{}_R[f]$ 
 (i.e., $\lambda_{\equiv_R}$) may cause $I$ to be invalid again,
 which needs another round of rebuilding. 
Therefore, the complete rebuilding function $R^*$ is defined as the Kleene star of $R$.

Finally, the evaluation of an \egglog program $P$ is defined as the Kleene star of $R^*\circ T_P$, 
 i.e., $[\![P]\!](I)=(R^*\circ T_P)^*(I)$. 
For many practice applications, 
 $[\![P]\!](I)$ often has infinite size, 
 so we only calcluated a finite approximation of result, 
 i.e., $\bigcup_{i=1,\ldots,n}(R^*\circ T_P)^n(I)$ for some iteration size $n$.

(Note: I think $[\![P]\!](I)=(R^*\circ T_P)^*(I)$ may
be more faithful to the implementation because 
in the implementation we never start with a fresh database
in each iteration. We accumulate)

(Note: Compared with Flix, we specifically does not constraint $T_P$ to be monotonic. 
In fact, \egglog programs can be non-monotonic. 
This is to be compatible with egg, as some egg applications use non-monotonic rules)

\newpage